% Part: many-valued-logic
% Chapter: syntax-and-semantics
% Section: introduction

\documentclass[../../../include/open-logic-section]{subfiles}

\begin{document}

\olfileid[fr]{mvl}{syn}{int}

\olsection{Introduction}

En logique classique, les formules
se construisent à partir de variables
propositionnelles au moyen des connecteurs
$\lnot$, $\land$, $\lor$, $\lif$ et
$\liff$. Sa sémantique emploie les deux
valeurs de vérité $\True$ et~$\False$.
Une valuation~$\pAssign{v}$ attribue
l'une de ces valeurs à chaque variable
propositionnelle. Elle détermine alors
une valeur~$\pValue{v}(!A)$ pour chaque
formule~$!A$. La formule~$!A$ est
satisfaite dans cette valuation,
$\pSat{v}{!A}$, si et seulement si
$\pValue{v}(!A)=\True$.

Les logiques plurivalentes généralisent
la logique classique en admettant
d'autres valeurs de vérité que
$\True$ et~$\False$. Une valuation
$\pAssign{v}$ attribue ainsi à chaque
variable propositionnelle~$p$ une
valeur prise dans un ensemble~$V$.
La logique classique en est un cas
où $V=\{\True,\False\}$; les tables
de vérité familières des connecteurs
permettent alors de calculer
$\pValue{v}(!A)$.

Dès que l'on ajoute d'autres valeurs,
plusieurs fonctions peuvent
interpréter les connecteurs que nous
lisons « et », « ou », « non » et
« si... alors ». Pour définir l'évaluation des formules, on fixe,
en plus de l'ensemble~$V$, une \emph{fonction de vérité} pour
chaque connecteur. Une fois ces fonctions fixées pour une
logique~$\Log L$, la valuation détermine de manière unique la
valeur $\pValue{v}(!A)[\Log L]$ de chaque formule. Comme la
logique classique, ces logiques sont \emph{vérifonctionnelles}.

Nous pouvons alors généraliser les notions
sémantiques de tautologie, de conséquence
et de satisfaisabilité. En logique
classique, une formule est une tautologie
si $\pValue{v}(!A)=\True$ pour toute
valuation~$\pAssign{v}$. Dans une
logique plurivalente, l'ensemble~$V$
ne contient pas nécessairement la
valeur~$\True$. Il peut, par exemple,
être formé de tous les rationnels
compris entre~$0$ et~$1$; la valeur~$1$
y joue généralement le rôle de~$\True$.
Dans d'autres logiques, plusieurs
valeurs jouent ce rôle. Nous distinguons
donc un sous-ensemble~$V^+\subseteq V$
de \emph{valeurs désignées}.

Une formule~$!A$ est satisfaite dans
la valuation~$\pAssign{v}$ pour la
logique~$\Log L$, ce que nous notons
$\pSat{v}{!A}[\Log L]$, si et seulement
si $\pValue{v}(!A)[\Log L]\in V^+$.
Elle est une tautologie de~$\Log L$,
$\Entails[\Log L] !A$, si cette
appartenance vaut pour toute valuation.
Enfin, un ensemble de formules~$\Gamma$
entraîne~$!A$, ce que nous notons
$\Gamma\Entails[\Log L] !A$, si toute
valuation qui satisfait toutes les
formules de~$\Gamma$ satisfait aussi~$!A$.

\end{document}
